\section{Training Methodology}
\label{sec:appendix_training}

This appendix provides the complete training specification for reproducibility.

\subsection{Common hyperparameters}

Table~\ref{tab:hyperparams} lists all hyperparameters shared across neural network models.

\begin{table}[htbp]
\centering
\caption{Training hyperparameters (all neural networks).}
\label{tab:hyperparams}
\begin{tabular}{ll}
\toprule
Hyperparameter & Value \\
\midrule
Optimizer                & Adam \citep{Kingma2015} \\
Initial learning rate    & $10^{-3}$ \\
Adam $\beta_1, \beta_2$  & 0.9, 0.999 \\
Adam $\epsilon$          & $10^{-8}$ \\
LR schedule              & Cosine annealing with warm restarts \\
Cosine $T_0$             & 200 epochs \\
Cosine $T_{\mathrm{mult}}$ & 2 \\
Cosine $\eta_{\min}$     & $10^{-6}$ \\
Batch size               & 256 \\
Max epochs               & 1{,}000 \\
Early stopping patience  & 150 epochs \\
Early stopping criterion & Validation MSE \\
Best model selection     & Lowest validation MSE across all epochs \\
Loss function            & MSE \\
Regularization           & None \\
Data precision (training)    & float32 \\
Data precision (features)    & float64 \\
\bottomrule
\end{tabular}
\end{table}

\subsection{Model architectures}

\begin{table}[htbp]
\centering
\caption{Neural network architecture details.}
\label{tab:architectures}
\begin{tabular}{lllrrl}
\toprule
Model & Layers & Activation & Parameters & Output & Target \\
\midrule
MLP       & $5 \to 32 \to 32 \to 1$         & Tanh / Linear & 1{,}249  & Scalar   & $|b|$ \\
MLP-Large & $5 \to 128 \to 128 \to 128 \to 128 \to 1$ & Tanh / Linear & 50{,}049 & Scalar   & $|b|$ \\
TBNN      & $5 \to 64 \to 64 \to 64 \to 10$ & Tanh / Linear & 9{,}354  & 10 coeff.\ & $b_{ij}$ \\
PI-TBNN   & Same as TBNN                     & Tanh / Linear & 9{,}354  & 10 coeff.\ & $b_{ij}$ \\
\bottomrule
\end{tabular}
\end{table}

\subsection{TBRF hyperparameters}

\begin{table}[htbp]
\centering
\caption{TBRF (scikit-learn \texttt{RandomForestRegressor}) hyperparameters.}
\label{tab:tbrf_hyperparams}
\begin{tabular}{ll}
\toprule
Parameter & Value \\
\midrule
Trees per coefficient    & 200 \\
Number of coefficients   & 10 (Pope basis) \\
Total trees              & 2{,}000 \\
Max depth                & 20 \\
Min samples split        & 2 (default) \\
Min samples leaf         & 1 (default) \\
Max features             & $\sqrt{n_{\mathrm{features}}}$ (default) \\
Bootstrap                & True \\
Random state             & 42 \\
\bottomrule
\end{tabular}
\end{table}

\subsection{Training results}

\begin{table}[htbp]
\centering
\caption{Training outcomes for all models.}
\label{tab:training_results}
\begin{tabular}{lrrrl}
\toprule
Model & Val RMSE & Epochs run & Stopped by & Training time \\
\midrule
TBRF (200 trees)         & 0.0637 & n/a   & n/a         & $\sim$500 s (4 CPU cores) \\
TBNN                     & 0.0845 & 694   & Early stop  & \todo{TBD} \\
PI-TBNN ($\beta=0.001$)  & 0.0852 & 729   & Early stop  & \todo{TBD} \\
PI-TBNN ($\beta=0.01$)   & 0.0909 & 676   & Early stop  & \todo{TBD} \\
MLP-Large                & 0.1045 & 344   & Early stop  & \todo{TBD} \\
MLP                      & 0.1096 & 1{,}000 & Max epochs & \todo{TBD} \\
\bottomrule
\end{tabular}
\end{table}

\subsection{Software and hardware}

\begin{table}[htbp]
\centering
\caption{Software versions used for training.}
\label{tab:software}
\begin{tabular}{ll}
\toprule
Package & Version \\
\midrule
Python     & 3.x (NVIDIA PyTorch container) \\
PyTorch    & 2.1.0a0+32f93b1 \\
NumPy      & 1.22.2 \\
scikit-learn & 1.2.0 \\
pandas     & 1.5.3 \\
\bottomrule
\end{tabular}
\end{table}

Training was performed on an NVIDIA L40S GPU at the Georgia Tech PACE cluster.
No fixed random seed was used for PyTorch (weight initialization and data shuffling are non-deterministic).
Feature extraction used float64 precision; training used float32.

\subsection{Weight export format}

Neural network weights are exported as human-readable text files with 10-digit scientific notation:
\begin{verbatim}
data/models/<model>_paper/
  layer0_W.txt    # [out_features, in_features]
  layer0_b.txt    # [out_features]
  layer1_W.txt
  ...
  input_means.txt # Z-score means [5]
  input_stds.txt  # Z-score stds [5]
  metadata.json   # Architecture info
\end{verbatim}

TBRF compact variants are exported as flat binary files with a 12-byte header (\texttt{total\_nodes}, \texttt{n\_basis}, \texttt{n\_trees} as int32) followed by five arrays (\texttt{children\_left}, \texttt{children\_right}, \texttt{feature}, \texttt{threshold}, \texttt{value}).

\section{Tensor Basis Functions}
\label{sec:appendix_basis}

The 10-tensor Pope integrity basis \citep{Pope1975} is computed from the non-dimensionalized strain rate $\Shat$ and rotation rate $\Omhat$:
\begin{align}
    T^{(1)} &= \Shat \\
    T^{(2)} &= \Shat\Omhat - \Omhat\Shat \\
    T^{(3)} &= \Shat^2 - \tfrac{1}{3}\mathrm{tr}(\Shat^2)I \\
    T^{(4)} &= \Omhat^2 - \tfrac{1}{3}\mathrm{tr}(\Omhat^2)I \\
    T^{(5)} &= \Omhat\Shat^2 - \Shat^2\Omhat \\
    T^{(6)} &= \Omhat^2\Shat + \Shat\Omhat^2 - \tfrac{2}{3}\mathrm{tr}(\Shat\Omhat^2)I \\
    T^{(7)} &= \Omhat\Shat\Omhat^2 - \Omhat^2\Shat\Omhat \\
    T^{(8)} &= \Shat\Omhat\Shat^2 - \Shat^2\Omhat\Shat \\
    T^{(9)} &= \Omhat^2\Shat^2 + \Shat^2\Omhat^2 - \tfrac{2}{3}\mathrm{tr}(\Shat^2\Omhat^2)I \\
    T^{(10)} &= \Omhat\Shat^2\Omhat^2 - \Omhat^2\Shat^2\Omhat
\end{align}
Each tensor is symmetric and traceless.
In the solver, only the 6 independent components $(11, 12, 13, 22, 23, 33)$ are stored and computed.

\section{Complete Profiling Tables}
\label{sec:appendix_profiling}

\todo{Add complete profiling tables for all GPU types (L40S, H100, H200), all flow cases, all models, including sub-phase breakdown.}
